\section{Implementation and Reproducibility}
\label{sec:impl}

\paragraph{Code organisation.}
The experiment pipeline is implemented under the project source directory with
modular components for configuration, data loading, autoregressive baselines,
speculative decoding, metric aggregation, and runtime orchestration.
The notebook driver executes these modules in fixed phase order and writes
artifacts to the results directory.

\paragraph{Model loading and precision.}
In the current configuration, target and draft models are loaded in FP16 for
the main sweep, with quantisation controls retained in configuration for
alternative runs. Tokenizer and generation limits are held constant across
baselines and speculative runs to preserve paired comparability.

\paragraph{Baselines.}
We run one autoregressive baseline per regime (deterministic and stochastic)
with the same manifests and output-length constraints as speculative runs.
Each baseline artifact stores per-sample latency, output length, and task
predictions.

\paragraph{Speculative decoding implementation.}
The speculative loop follows standard token-level draft--verify logic.
At each cycle, the draft proposes up to $k$ tokens autoregressively; the
target verifies the proposal in one pass over the extended prefix; the accepted
prefix is committed up to the first mismatch; and one correction token is added
from the target distribution when needed. This procedure preserves target-level
generation semantics while exposing acceptance statistics needed for analysis.
Figure~\ref{fig:specdec-impl-flow} summarizes this implementation flow.

\begin{figure}[t]
\centering
\resizebox{0.98\columnwidth}{!}{%
\begin{tikzpicture}[
	node distance=4.8mm and 6mm,
	box/.style={rectangle, draw, rounded corners=1mm, align=center, minimum height=6mm, text width=31mm, inner sep=1.5mm},
	decision/.style={diamond, draw, align=center, aspect=2, inner sep=1.1mm, text width=22mm},
	line/.style={-Latex, thick, shorten >=1pt, shorten <=1pt}
]
\node[box] (start) {Start sample\\encode prompt and init KV cache};
\node[box, below=of start] (draft) {Draft model proposes\\up to $k$ tokens};
\node[box, below=of draft] (verify) {Target verifies draft block\\in one forward pass};
\node[decision, below=of verify] (accept) {All $k$ tokens\\accepted?};
\node[box, left=of accept, xshift=-9mm] (bonus) {Commit $k$ tokens\\+ one target bonus token};
\node[box, right=of accept, xshift=9mm] (correct) {Commit accepted prefix\\+ one corrected target token};
\node[box, below=of accept, yshift=-2mm] (crop) {Crop/update draft and target\\KV caches to committed prefix};
\node[decision, below=of crop] (stop) {EOS or max tokens\\reached?};
\node[box, below=of stop] (end) {Finish sample and\\record runtime/quality fields};

\draw[line] (start) -- (draft);
\draw[line] (draft) -- (verify);
\draw[line] (verify) -- (accept);
\draw[line] (accept) -- node[above, font=\scriptsize] {yes} (bonus);
\draw[line] (accept) -- node[above, font=\scriptsize] {no} (correct);
\draw[line] (bonus) |- (crop.west);
\draw[line] (correct) |- (crop.east);
\draw[line] (crop) -- (stop);
\draw[line] (stop) -- node[right, font=\scriptsize] {yes} (end);
\draw[line] (stop.west) -- ++(-14mm,0) |- node[pos=0.25, left, font=\scriptsize] {no} (draft.west);
\end{tikzpicture}%
}
\caption{Vanilla speculative decoding implementation flow used in Phases 3--4.
The target model remains the verifier of record; accepted draft prefixes are
committed and caches are cropped to preserve exact target-consistent decoding.}
\label{fig:specdec-impl-flow}
\end{figure}

\paragraph{Regimes and outputs.}
The same verification core is used for deterministic and stochastic modes,
with only sampling policy changed by regime parameters.
Each configuration writes a dedicated CSV artifact with per-sample runtime,
token counts, and quality-relevant outputs.

\paragraph{Runtime cost model for CUDA-graph path.}
For a verify cycle, we decompose runtime into launch/scheduling and kernel
compute components. In eager execution,
\begin{equation}
T_{\text{eager}} \approx N\,(t_{\text{launch}} + t_{\text{kernel}} + t_{\text{sync}}).
\end{equation}
With CUDA-graph replay ratio $\rho$ and one-time capture cost
$t_{\text{capture}}$, runtime becomes
\begin{equation}
T_{\text{graph}} \approx t_{\text{capture}} + N\,[\rho\,(t_{\text{replay}} + t_{\text{kernel}})
+ (1-\rho)(t_{\text{launch}} + t_{\text{kernel}} + t_{\text{sync}})].
\end{equation}
Graph execution helps when replay is stable ($\rho \rightarrow 1$) and
$t_{\text{replay}} \ll t_{\text{launch}}$. On short generation windows,
this can offset speculative verification overhead by reducing repeated host
launch costs.

In our current run family, graph capture metadata remains pending with zero
capture rate, so the main reported gains are from the optimized code path
under eager fallback. An auxiliary short-sequence toggle benchmark
($n{=}12$, $k{=}8$, 64 output tokens) shows +3.37\% TTFT, +1.41\% derived TPOT,
and +1.56\% end-to-end latency for \texttt{gpu\_opt\_on} versus
\texttt{gpu\_opt\_off}, consistent with the expectation that capture instability
reduces the benefit of launch-optimization techniques.
